\documentclass[11pt]{article}
\usepackage[T1]{fontenc}
\usepackage{textcomp}
\usepackage[margin=0.75in]{geometry}
\usepackage{amsmath,amssymb,amsthm}
\usepackage{array,booktabs}
\usepackage{hyperref}
\setlength{\parindent}{0pt}
\newtheorem{theorem}{Theorem}
\theoremstyle{definition}
\newtheorem{definition}{Definition}
\title{Flow Proof Artifact: Eq}
\author{Generated by \texttt{flow doc proof}}
\date{Flow Proof Book}
\begin{document}
\maketitle
Leibniz identity — equality reflexivity on natural numbers.\\[0.5em]
\textbf{Source.} leibniz — https://en.wikipedia.org/wiki/Law\_of\_identity\\[0.5em]
\bigskip
\noindent{\large \textbf{Axiom 1.} \textit{Anything equals itself} \hfill equality \cdot equality \cdot everything equals itself}\par
\smallskip\noindent\textit{Coordinate: equality \cdot equality \cdot everything equals itself}\par
\smallskip\noindent\textit{Source: leibniz — https://en.wikipedia.org/wiki/Law\_of\_identity}\par
\medskip
\noindent\textbf{Goal.} Anything equals itself.  (42 == 42)\par
\noindent$\displaystyle \forall x \in \mathbb{N}\quad x = x$\par
\medskip
\noindent
\begin{tabular}{@{} >{\raggedright\arraybackslash}p{0.06\textwidth} >{\raggedright\arraybackslash}p{0.47\textwidth} >{\raggedleft\arraybackslash}p{0.41\textwidth} @{}}
\textbf{\#} & \textbf{Proof} & \textbf{Mathematics} \\
\midrule
\textbf{1.} & We accept everything equals itself for equality on equality without proof — an ontological commitment, not a lemma. & \\[0.45em]
\textbf{2.} & This holds immediately by the stated axiom: x equals x. Hence proven. & $\displaystyle x = x$ \\[0.45em]
\bottomrule
\end{tabular}
\label{thm:Eq-equality-everything_equals_itself}
\par\medskip\hrule
\end{document}
